% ---------------------------------------------------------------------------
%  Schemas TikZ du dossier T005
%  Unite : le millimetre. La largeur utile du texte vaut 160 mm.
% ---------------------------------------------------------------------------

% Ligne secondaire d'un noeud : la fonte doit etre reappliquee ligne par ligne,
% car l'alignement d'un noeud TikZ place chaque ligne dans son propre groupe.
\newcommand{\fline}[1]{{\tiny\ttfamily #1}}

% ------------------------------------------------------- figure : architecture
\newcommand{\figArchitecture}{\resizebox{\textwidth}{!}{%
\begin{tikzpicture}[x=1mm,y=1mm,every node/.style={font=\scriptsize}]

  % ---- panneau 1 : repartition physique
  \node[titrepan] at (0,93) {Panneau 1 \quad Répartition des quatre instances
                             sur les trois hôtes};

  \foreach \x/\nom/\mgmt in {0/H1/10.30.0.11, 54/H2/10.30.0.12, 108/H3/10.30.0.13}{
    \draw[hote] (\x,58) rectangle ++(50,30);
    \node[font=\scriptsize\bfseries,text=gdnavy,anchor=north west] at (\x+2,87) {Hôte~\nom};
    \node[font=\tiny\ttfamily,text=gdgrey,anchor=north east] at (\x+48,86.5) {\mgmt};
  }

  \node[blocteal,minimum width=45mm,minimum height=9mm,anchor=south west] at (2.5,71)
    {\textbf{ns1} \,\ autoritatif primaire\\[-2pt]{\tiny\ttfamily 10.30.4.10}};
  \node[blocteal,minimum width=45mm,minimum height=9mm,anchor=south west] at (2.5,60)
    {\textbf{rec2} \,\ résolveur récursif\\[-2pt]{\tiny\ttfamily 10.30.4.13}};
  \node[blocteal,minimum width=45mm,minimum height=9mm,anchor=south west] at (56.5,71)
    {\textbf{rec1} \,\ résolveur récursif\\[-2pt]{\tiny\ttfamily 10.30.4.12}};
  \node[blocteal,minimum width=45mm,minimum height=9mm,anchor=south west] at (110.5,71)
    {\textbf{ns2} \,\ autoritatif secondaire\\[-2pt]{\tiny\ttfamily 10.30.4.11}};

  \node[font=\tiny,text=gdgrey,anchor=north west,text width=158mm] at (0,56)
    {Anti-affinité par rôle : les deux membres d'une même paire ne résident jamais
     sur le même hôte, de sorte que toute perte d'un hôte laisse un serveur de
     chaque rôle en service.};

  \draw[gdrule,line width=0.4pt] (0,49) -- (158,49);

  % ---- panneau 2 : chemins logiques
  \node[titrepan] at (0,45) {Panneau 2 \quad Chemins de résolution};

  % chemin direct interdit, place dans un couloir degage au-dessus des blocs
  \node[etiqr,font=\tiny\bfseries] at (79,39.8)
    {accès direct d'un VNet tenant aux serveurs autoritatifs : refusé au pare-feu};
  \draw[draw=gdred,line width=0.65pt,dash pattern=on 2.4pt off 1.8pt]
       (20,32) -- (20,36.4) -- (140,36.4) -- (140,32);
  \draw[gdred,line width=0.9pt] (77.8,35.2) -- (80.2,37.6);
  \draw[gdred,line width=0.9pt] (77.8,37.6) -- (80.2,35.2);

  \node[bloc,minimum width=40mm,minimum height=16mm,anchor=south west] (clients) at (0,16)
    {\textbf{Clients autorisés}\\[1pt]
     \fline{VNets tenant 10.30.64.0/18}\\
     \fline{gestion 10.30.0.0/24}\\
     \fline{VPN admin 10.30.16.32/28}};
  \node[blocteal,minimum width=40mm,minimum height=16mm,anchor=south west] (rec) at (60,16)
    {\textbf{Résolveurs rec1 et rec2}\\[1pt]
     \fline{10.30.4.12 et .13}\\
     \fline{allow-from puis vue}\\
     \fline{cache indexé par étiquette}};
  \node[bloc,minimum width=40mm,minimum height=16mm,anchor=south west] (auth) at (120,16)
    {\textbf{Autoritatifs ns1 et ns2}\\[1pt]
     \fline{10.30.4.10 et .11}\\
     \fline{zone gandal.internal}\\
     \fline{et 30.10.in-addr.arpa}};

  \draw[fl] (clients) -- node[etiqt,above=1pt,font=\tiny\ttfamily]{UDP/TCP 53} (rec);
  \draw[flnavy] (rec) -- node[etiq,above=1pt,font=\tiny\ttfamily]{forward-zones} (auth);

  \node[blocgris,minimum width=40mm,minimum height=12mm,anchor=south west] (racine) at (60,0)
    {\textbf{Racine publique}\\[1pt]
     \fline{via PF-WAN puis R1}\\
     \fline{dnssec=validate}};
  \node[blocgold,minimum width=40mm,minimum height=12mm,anchor=south west] (repl) at (120,0)
    {\textbf{Réplication}\\[1pt]
     \fline{NOTIFY puis AXFR ou IXFR}\\
     \fline{signés TSIG, ns1 vers ns2}};

  \draw[flgris] (rec) -- (racine);
  \draw[flgold] (auth) -- (repl);

  \node[font=\tiny,text=gdgrey,anchor=north west,text width=158mm] at (0,-4)
    {La carte réseau unique par hôte et le commutateur SW1 unique restent des
     dépendances communes. La redondance couvre la perte d'une machine virtuelle
     ou d'un hôte, et non celle de SW1 ou du routeur R1.};
\end{tikzpicture}}}

% -------------------------------------------------------------- figure : zones
\newcommand{\figZones}{\resizebox{\textwidth}{!}{%
\begin{tikzpicture}[x=1mm,y=1mm,every node/.style={font=\scriptsize}]

  \node[titrepan] at (0,62) {Espace direct};
  \node[titrepan] at (84,62) {Espace inverse};

  \node[bloc,minimum width=74mm,minimum height=12mm,anchor=north west] (gi) at (0,58)
    {\textbf{gandal.internal}\\[1pt]
     \fline{SOA, NS ns1 et ns2, adresses des 4 serveurs}\\
     \fline{délégations NS vers les zones filles}};
  \node[blocteal,minimum width=66mm,minimum height=10mm,anchor=north west] (gia) at (8,41)
    {\textbf{infra.gandal.internal}\\[1pt]
     \fline{hôtes, XO, bastion, pfSense, services}};
  \node[blocteal,minimum width=66mm,minimum height=10mm,anchor=north west] (gib) at (8,27)
    {\textbf{tenant-001 à tenant-050}\\[1pt]
     \fline{50 zones, une par tenant}};
  \draw[gdgrey,line width=0.4pt] (3.5,46) |- (gia.west);
  \draw[gdgrey,line width=0.4pt] (3.5,46) |- (gib.west);

  \node[bloc,minimum width=74mm,minimum height=12mm,anchor=north west] (ri) at (84,58)
    {\textbf{30.10.in-addr.arpa}\\[1pt]
     \fline{le /16 appartient entièrement au projet}\\
     \fline{aucun découpage RFC 2317 nécessaire}};
  \node[blocteal,minimum width=66mm,minimum height=10mm,anchor=north west] (ria) at (92,41)
    {\textbf{11 zones d'infrastructure}\\[1pt]
     \fline{octets 0, 1, 2, 4, 5, 6, 7, 8, 10, 16 et 17}};
  \node[blocteal,minimum width=66mm,minimum height=10mm,anchor=north west] (rib) at (92,27)
    {\textbf{50 zones tenant}\\[1pt]
     \fline{octets 64 à 113, une zone par /24}};
  \draw[gdgrey,line width=0.4pt] (87.5,46) |- (ria.west);
  \draw[gdgrey,line width=0.4pt] (87.5,46) |- (rib.west);

  \node[bande,minimum width=155mm,anchor=north west,text width=150mm,align=left] at (0,12)
    {{\tiny\bfseries\scshape\color{gdteal} Correspondance entre un tenant, sa zone
      directe et sa zone inverse}\\[3pt]
     $\mathcal{D}(n)=\texttt{tenant-0}n\texttt{.gandal.internal}$ \quad et \quad
     $\mathcal{R}(n)=(63+n)\texttt{.30.10.in-addr.arpa}$, \quad
     pour $1\leqslant n\leqslant 50$\\[3pt]
     {\tiny Exemple pour $n=7$ : \texttt{tenant-007.gandal.internal} et
      \texttt{70.30.10.in-addr.arpa}, préfixe \texttt{10.30.70.0/24}.}};
\end{tikzpicture}}}

% ----------------------------------------------------------- figure : sequence
\newcommand{\figSequence}{\resizebox{\textwidth}{!}{%
\begin{tikzpicture}[x=1mm,y=1mm,every node/.style={font=\scriptsize}]

  \foreach \x/\y/\n/\ti/\la/\lb in {%
     0/44/1/{Bail attribué}/{Kea délivre 10.30.70.37}/{à la VIF de la machine},
    56/44/2/{Publication du A}/{agent vers l'API de ns1}/{PATCH rrset, type REPLACE},
   112/44/3/{Publication du PTR}/{zone 70.30.10.in-addr.arpa}/{même operation\_id},
     0/20/4/{Réplication}/{NOTIFY de ns1 vers ns2}/{AXFR ou IXFR signé TSIG},
    56/20/5/{Contrôle}/{requête sur ns1, ns2,}/{rec1 et rec2 ; concordance},
   112/20/6/{État READY}/{exposé au Backend}/{operation\_id conservé}}{%
    \draw[blocteal,fill=white] (\x,\y) rectangle ++(46,17);
    \fill[gdteal] (\x,\y+13) rectangle ++(46,4);
    \node[font=\tiny\bfseries,text=white,anchor=west]  at (\x+1.6,\y+15) {\n};
    \node[font=\tiny\bfseries,text=white]              at (\x+25,\y+15) {\ti};
    \node[font=\tiny\ttfamily,text=gdgrey]             at (\x+23,\y+8.5) {\la};
    \node[font=\tiny\ttfamily,text=gdgrey]             at (\x+23,\y+4.5) {\lb};
  }

  \draw[flnavy] (46,50) -- (56,50);
  \draw[flnavy] (102,50) -- (112,50);
  \draw[flnavy] (46,26) -- (56,26);
  \draw[flnavy] (102,26) -- (112,26);
  \draw[flnavy] (135,44) -- (135,40) -- (23,40) -- (23,37);

  \draw[flrouge] (79,20) -- (79,15);
  \draw[blocrouge,fill=gdwarn] (14,4) rectangle ++(130,11);
  \node[font=\tiny\bfseries,text=gdred] at (79,11.5)
    {Échec de publication ou de contrôle : l'opération reste observable en état dégradé,};
  \node[font=\tiny,text=gdred] at (79,7)
    {l'adresse n'est pas rendue au pool et une alerte est levée.};
\end{tikzpicture}}}

% ------------------------------------------------------------- figure : chemin
\newcommand{\figChemin}{\resizebox{\textwidth}{!}{%
\begin{tikzpicture}[x=1mm,y=1mm,every node/.style={font=\scriptsize}]

  \node[bloc,minimum width=34mm,minimum height=13mm,anchor=south west] (vm) at (0,64)
    {\textbf{VM du tenant 007}\\[1pt]
     \fline{10.30.70.37}\\
     \fline{resolv.conf : .12 puis .13}};
  \node[bloc,minimum width=34mm,minimum height=13mm,anchor=south west] (gw) at (41,64)
    {\textbf{Passerelle .1}\\[1pt]
     \fline{paire PF-USER}\\
     \fline{VIP 10.30.70.1}};
  \node[blocrouge,minimum width=34mm,minimum height=13mm,anchor=south west] (pf) at (82,64)
    {\textbf{Filtre pfSense}\\[1pt]
     \fline{53 vers .12 et .13}\\
     \fline{.10 et .11 refusés}};
  \node[blocteal,minimum width=34mm,minimum height=13mm,anchor=south west] (re) at (123,64)
    {\textbf{rec1 ou rec2}\\[1pt]
     \fline{allow-from vérifié}\\
     \fline{gettag() calcule la vue}};
  \draw[fl] (vm) -- (gw);   \draw[fl] (gw) -- (pf);   \draw[fl] (pf) -- (re);

  \node[bande,minimum width=155mm,anchor=north west,text width=150mm,align=left] at (0,58)
    {{\tiny\bfseries\scshape\color{gdteal} preresolve() : politique de sélection
      des zones}\\[3pt]
     {\tiny\ttfamily nom sous tenant-007.gandal.internal ou 70.30.10.in-addr.arpa
      \ :\ résolution autorisée}\\[1.5pt]
     {\tiny\ttfamily\color{gdred} nom sous infra.gandal.internal ou sous un autre
      tenant \ :\ REFUSED, sans divulgation}\\[1.5pt]
     {\tiny\ttfamily\color{gdgrey} nom public \ :\ récursion ordinaire, quelle que
      soit la vue}};
  \draw[fl] (140,64) -- (140,58);

  \node[bloc,minimum width=46mm,minimum height=13mm,anchor=south west] (na) at (0,14)
    {\textbf{Autoritatifs ns1 et ns2}\\[1pt]
     \fline{10.30.4.10 et .11}\\
     \fline{atteints par forward-zones}};
  \node[blocgris,minimum width=46mm,minimum height=13mm,anchor=south west] (rp) at (56,14)
    {\textbf{Racine publique}\\[1pt]
     \fline{via PF-WAN puis R1}\\
     \fline{dnssec=validate}};
  \node[blocgold,minimum width=46mm,minimum height=13mm,anchor=south west] (jo) at (112,14)
    {\textbf{Journal dnstap}\\[1pt]
     \fline{horodatage en UTC}\\
     \fline{lu en pull par Monitoring}};
  \draw[flnavy] (23,34) -- (23,27);
  \draw[flgris] (79,34) -- (79,27);
  \draw[flgold] (135,34) -- (135,27);

  \node[font=\tiny,text=gdgrey,anchor=north west,text width=158mm] at (0,12)
    {\textbf{\color{gdteal}L'étiquette de vue est calculée avant toute consultation
     de cache et entre dans la clé du cache paquet : deux vues distinctes ne peuvent
     jamais partager une réponse.} Un contrôle placé dans \texttt{preresolve()} seul
     n'offrirait pas cette garantie.};
\end{tikzpicture}}}
